\documentclass[五凉全誌六德集]{ltc-guji}
\无标点模式
\setmainfont{TW-Kai}

\title{武威縣誌}
\chapter{卷一}
\出版商{智集}

\begin{document}
\begin{正文}
五凉全誌六德集前序
\begin{段落}[indent=2]
嘗思以道事\平抬 君\圈点{者，敬事爲本，而文獻無徵，施之或謬，則學術之限也}。夫述經論\圈点{史}，探立治之源\圈点{流，而在此一地與此一時，宜如何措置，則又不可泛涉概擬，而不求其當也。夫誌書者，一方之文獻也。}前之所歷，後之所鑑，即或人情運會，地與時易，而轉移不甚相遠。考諸巳徃，以酌今兹，如子孫之則效乎。高曾祖父無慮迂迴，岐黄之騐藥，於既試無慮悞投，切莫切於此矣。然而誌書之爲乙類頗繁，非從簡易，則目眩神離，奚能據以斟酌？五凉舊誌缺略，謀所以輯之，未得其要。永昌司鐸上官鈍齊示以對山武功誌因以爲式，定爲七萹。一地里誌，星野沿革、疆域圖，並説山川、村社、户口、田畝、賦則、物産、水利圖，並説古蹟、祥異寓焉。天時、地利渾而爲一，是飬民之原也。一建置誌，城廓、公署、壇壝、學校、驛傳、寺觀寓焉。規模宏整，報本崇文，是教民之由也。一風俗誌，四民執業，五禮彙儀，異端異族寓焉，是\平抬 聖天子大同之模，齊民之具也。一官師誌，職官、名宦寓焉，表制度紀功德，是重民之司牧也。一官師誌，職官、名宦寓焉，表制度紀功德，是重民之司牧也。一兵防誌營堡、關隘、烽墩、軍制，糧餉寓焉，安不忘危，是邊陰之所以衛民也。一人物誌鄉賢、忠孝、節義、選舉、流寓寓焉。揭懿美於前徽，是鼓舞斯民之機也。一文藝誌議疏、碑記、詩歌寓焉，君子學道則愛人，小人學道，則易使也。是嫻民於禮樂也。立治者於此，展卷了然，即其類以求之，自必有得乎最當者。以是爲敬事，而非泛涉者比也；以是爲以道事，\平抬 君而非文獻無徴者比也。寅寮其勖諸！愚於五凉將及五載，無日不以兢業爲懷，以冀上不負。\平抬 君下不累民，而識淺才疏，終歸惕厲。
所著有敬時考績十萹，天山自訓、訓吏、訓故、士訓民十萹，勸課歌三十二萹，勸農示十二萹，以及偶與同官縉紳、耆老兒童所口占而筆繪者，集爲一册曰學道編，附於誌末，俾俊之覽者以爲不徒記載之工，而几我士民永奉爲
\end{段落}

\end{正文}
\end{document}